\documentclass[a4paper,11pt]{article}

% Layout
\usepackage[a4paper, top=2.2cm, bottom=2.2cm, left=2.5cm, right=2.5cm]{geometry}

% Encoding / font
\usepackage[T1]{fontenc}
\usepackage[utf8]{inputenc}
\usepackage{lmodern}
\usepackage{microtype}

% Colour
\usepackage{xcolor}
\definecolor{headblue}{HTML}{1B3A6B}
\definecolor{codebg}{HTML}{F4F4F4}
\definecolor{warnbg}{HTML}{FFF3E0}
\definecolor{infobg}{HTML}{E8F4FD}

% Coloured callout boxes
\usepackage{tcolorbox}
\newtcolorbox{warnbox}{
  colback=warnbg,
  colframe=orange!70!red,
  boxrule=1pt,
  left=8pt, right=8pt, top=5pt, bottom=5pt,
  arc=2pt,
}
\newtcolorbox{infobox}{
  colback=infobg,
  colframe=blue!55!black,
  boxrule=1pt,
  left=8pt, right=8pt, top=5pt, bottom=5pt,
  arc=2pt,
}

% Code listings
\usepackage{listings}
\lstset{
  basicstyle=\ttfamily\small,
  backgroundcolor=\color{codebg},
  breaklines=true,
  breakatwhitespace=false,
  postbreak=\mbox{\textcolor{gray}{$\hookrightarrow$}\ },
  frame=single,
  framerule=0.4pt,
  rulecolor=\color{gray!50},
  xleftmargin=4pt,
  xrightmargin=4pt,
  aboveskip=4pt,
  belowskip=4pt,
  columns=flexible,
  keepspaces=true,
  showstringspaces=false,
}

% Section headings
\usepackage{titlesec}
\titleformat{\section}
  {\large\bfseries\color{headblue}}{}{0em}{}
  [\vspace{1pt}\color{headblue!60}\hrule\vspace{2pt}]
\titleformat{\subsection}
  {\normalsize\bfseries\color{headblue!80}}{}{0em}{}
\titlespacing{\section}{0pt}{14pt}{5pt}
\titlespacing{\subsection}{0pt}{9pt}{3pt}

% Paragraph spacing
\usepackage{parskip}
\setlength{\parskip}{5pt}

% Lists
\usepackage{enumitem}
\setlist[enumerate]{topsep=2pt, itemsep=2pt, parsep=0pt, leftmargin=*}
\setlist[itemize]{topsep=2pt, itemsep=2pt, parsep=0pt, leftmargin=*}

% Tables
\usepackage{tabularx}
\usepackage{booktabs}
\usepackage{array}

% Header / footer
\usepackage{fancyhdr}
\pagestyle{fancy}
\fancyhf{}
\lhead{\small\color{gray!70} DashyDashboard --- Production Deployment Guide}
\rhead{\small\color{gray!70} Broadridge Internal}
\cfoot{\small\thepage}
\renewcommand{\headrulewidth}{0.4pt}

% Links
\usepackage[hidelinks]{hyperref}

% ================================================================
\begin{document}
% ================================================================

\thispagestyle{empty}

\begin{center}
  {\LARGE\bfseries\color{headblue} DashyDashboard}\\[5pt]
  {\large\bfseries Production Deployment Guide}\\[5pt]
  {\small ASP.NET 7 \textbf{(net7.0)} \quad|\quad
         IIS \quad|\quad
         Windows Authentication \quad|\quad
         SQL Server}\\[3pt]
  {\small Target SQL Server: \texttt{clipvwbpod02} \quad|\quad
         Database: \texttt{ClientsAppAttestation}}\\[3pt]
  {\small\color{gray} June 2026}
\end{center}

\vspace{4pt}
\noindent{\color{headblue}\rule{\linewidth}{1.2pt}}
\vspace{6pt}

\tableofcontents
\newpage

% ================================================================
\section{Do This First}
% ================================================================

Work through these steps in order. Each step has a section below that tells you exactly what to do.

\begin{enumerate}
  \item \textbf{Read Section~\ref{sec:blockers} right now.}
        There are known database problems that will corrupt data if you go live too early.
  \item Install the right software on the build machine and on the IIS server
        (Section~\ref{sec:prereqs}).
  \item Edit the connection string file in the source repository,
        \emph{before} you publish (Sections~\ref{sec:whichfile} and~\ref{sec:whatpaste}).
  \item Run the one-line publish command on the build machine
        (Section~\ref{sec:publish}, Step~3).
  \item Copy the published output folder to the IIS server
        (Section~\ref{sec:publish}, Step~4).
  \item Run database migrations from a machine that can reach \texttt{clipvwbpod02}
        (Section~\ref{sec:publish}, Step~5).
  \item Create the IIS application pool and website, enable Windows Authentication,
        disable Anonymous Authentication (Section~\ref{sec:winauth}).
  \item Confirm all three \texttt{DeveloperMode} flags are \texttt{false}
        (Section~\ref{sec:demodata}).
\end{enumerate}

\begin{warnbox}
\textbf{STOP before you point this at the live database.}
Read Section~\ref{sec:blockers} first.
Known schema mismatches in the current code will cause data corruption and leave
Admin/GFH/IFH role enforcement completely bypassed.
\end{warnbox}

% ================================================================
\section{What to Install and Where}
\label{sec:prereqs}
% ================================================================

Two machines are involved.
The \textbf{build machine} is where you compile and publish the code.
The \textbf{IIS server} is the Windows Server that users browse to.
They need different software.

\bigskip
\renewcommand{\arraystretch}{1.6}
\begin{tabularx}{\linewidth}{|>{\bfseries}p{4.5cm}|X|X|}
\hline
Software & Build machine & IIS server \\
\hline
.NET SDK 7.0.x or later &
  \textbf{Required.}
  The app targets \texttt{net7.0}. SDK 7.0.x is the minimum.
  SDK 8.0.x also builds net7.0 projects and is acceptable. &
  Not needed. \\
\hline
.NET 7 Hosting Bundle &
  Not needed. &
  \textbf{Required.}
  Installs the ASP.NET Core 7 runtime and the IIS module that runs the app. \\
\hline
Node.js 18 or later &
  \textbf{Required.}
  The publish command runs \texttt{npm ci} and \texttt{npm run build} automatically. &
  \textbf{Not needed at all.}
  React is compiled into static files during publish. \\
\hline
IIS &
  Not needed. &
  \textbf{Required.} \\
\hline
\end{tabularx}
\renewcommand{\arraystretch}{1}

\bigskip
\begin{infobox}
\textbf{SDK vs Runtime.}
The project file targets \texttt{net7.0}.
The \textbf{server} only ever needs the .NET~7 runtime (the Hosting Bundle).
The \textbf{build machine} needs the .NET~7 SDK (minimum);
SDK~8.x also builds net7.0 projects so either works.
\textbf{Important:} .NET~7 is out of Microsoft mainstream support.
Plan an upgrade to .NET~8 or .NET~9 in a future sprint.
\end{infobox}

\subsection*{Install .NET SDK 7.0.x on the build machine}

\begin{enumerate}
  \item Open a browser. Go to \texttt{https://dotnet.microsoft.com/en-us/download/dotnet/7.0}
  \item Under \textbf{SDK}, click the \textbf{Windows x64 Installer} link. Download the file.
  \item Double-click the downloaded \texttt{.exe} file.
        A Windows security prompt appears. Click \textbf{Yes}.
  \item Click \textbf{Install}. Wait for the progress bar to finish. Click \textbf{Close}.
  \item Open a \textbf{new} Command Prompt window. (Press \textbf{Win + R},
        type \texttt{cmd}, press \textbf{Enter}.) Run:
\begin{lstlisting}
dotnet --version
\end{lstlisting}
        The output must begin with \texttt{7.} (or \texttt{8.} if you have SDK 8 installed --- both work).
\end{enumerate}

\subsection*{Install Node.js 18 LTS on the build machine}

\begin{enumerate}
  \item Open a browser. Go to \texttt{https://nodejs.org/en/download}
  \item Click the \textbf{Windows Installer (.msi)} link for version \textbf{18.x LTS}.
        Download the file.
  \item Double-click the \texttt{.msi} file. Click \textbf{Next} through every screen.
        Accept the licence. Leave all options at their defaults.
        Click \textbf{Install}. Accept the UAC prompt. Click \textbf{Finish}.
  \item Open a \textbf{new} Command Prompt window. Run:
\begin{lstlisting}
node --version
\end{lstlisting}
        The output must begin with \texttt{v18.}
\end{enumerate}

\subsection*{Install the .NET 7 Hosting Bundle on the IIS server}

Do this on the \textbf{IIS server}, not on the build machine.
Log in to the IIS server directly or via Remote Desktop.

\begin{enumerate}
  \item Open a browser on the IIS server.
        Go to \texttt{https://dotnet.microsoft.com/en-us/download/dotnet/7.0}
  \item Under \textbf{ASP.NET Core Runtime 7.0.x}, click the \textbf{Hosting Bundle} link.
        Download the file.
  \item Double-click the installer. Click \textbf{Install}. Accept the UAC prompt.
        Wait for it to finish. Click \textbf{Close}.
  \item Open a Command Prompt \textbf{as Administrator}:
        click \textbf{Start}, type \texttt{cmd},
        right-click \textbf{Command Prompt}, click \textbf{Run as administrator}. Run:
\begin{lstlisting}
dotnet --list-runtimes
\end{lstlisting}
        You should see a line containing \texttt{Microsoft.AspNetCore.App 7.0.}
  \item Still as Administrator, run:
\begin{lstlisting}
iisreset
\end{lstlisting}
        IIS must be restarted so it loads the new ASP.NET Core module.
\end{enumerate}

% ================================================================
\section{Exactly Which File Do I Edit?}
\label{sec:whichfile}
% ================================================================

\subsection*{Before publishing (working inside the source repository)}

Open this file in any text editor --- Notepad, Notepad++, or VS Code.
The path below is relative to the root of the repository:

\begin{lstlisting}
DashyDashboard/DashyDashboard.Api/appsettings.Production.json
\end{lstlisting}

Edit this file \emph{before} you run the publish command.
The publish step copies this file into the output folder automatically.

\subsection*{After publishing (on the IIS server)}

If you have already published and just need to update the connection string
without republishing, edit the file at the physical path of the IIS site:

\begin{lstlisting}
<IIS site physical path>\appsettings.Production.json
\end{lstlisting}

For example, if your IIS site's physical path is
\texttt{C:\textbackslash inetpub\textbackslash DashyDashboard}, the file you edit is:

\begin{lstlisting}
C:\inetpub\DashyDashboard\appsettings.Production.json
\end{lstlisting}

To open it:
open File Explorer, navigate to \texttt{C:\textbackslash inetpub\textbackslash DashyDashboard},
right-click \texttt{appsettings.Production.json}, click \textbf{Open with}, click
\textbf{Notepad}.

\subsection*{How the app decides which settings file to load}

The app loads files in this order.
Values in later entries override earlier ones.

\begin{enumerate}
  \item \texttt{appsettings.json} --- base file, always loaded, contains no connection string
  \item \texttt{appsettings.Production.json} --- loaded only when
        \texttt{ASPNETCORE\_ENVIRONMENT=Production}; contains the connection string
  \item Environment variables --- override both files;
        can be set in IIS Manager or in \texttt{web.config}
\end{enumerate}

If \texttt{appsettings.Production.json} is missing or the connection string is blank,
the app throws this error at startup and refuses to start:
\emph{``Connection string `Default' is not configured. Set it in
appsettings.Production.json or via environment variables.''}

\subsection*{Tell the app it is running in Production}

The app must know it is in the Production environment or it will not load the
\texttt{appsettings.Production.json} file.
Set the \texttt{ASPNETCORE\_ENVIRONMENT} variable in the published \texttt{web.config}.

\begin{enumerate}
  \item On the IIS server, open File Explorer.
        Navigate to the site folder, e.g.\ \texttt{C:\textbackslash inetpub\textbackslash DashyDashboard}.
  \item Right-click \texttt{web.config}. Click \textbf{Open with}. Click \textbf{Notepad}.
  \item Find the line that starts with \texttt{<aspNetCore}.
        Look for an \texttt{<environmentVariables>} block inside it.
        If that block does not exist, add it.
        The result should look exactly like this:
\begin{lstlisting}[language={}]
<aspNetCore processPath=".\DashyDashboard.Api.exe"
            stdoutLogEnabled="false"
            stdoutLogFile=".\logs\stdout"
            hostingModel="inprocess">
  <environmentVariables>
    <environmentVariable
        name="ASPNETCORE_ENVIRONMENT"
        value="Production" />
  </environmentVariables>
</aspNetCore>
\end{lstlisting}
  \item Press \textbf{Ctrl + S} to save. Close Notepad.
        IIS picks up \texttt{web.config} changes immediately --- no restart needed.
\end{enumerate}

% ================================================================
\section{What Do I Paste Into It?}
\label{sec:whatpaste}
% ================================================================

Open \texttt{appsettings.Production.json} in Notepad.
Press \textbf{Ctrl + A} to select everything.
Press \textbf{Delete} to clear the file.
Now paste this block exactly as it appears below:

\begin{lstlisting}[language={}]
{
  "ConnectionStrings": {
    "Default": "Server=clipvwbpod02;Database=ClientsAppAttestation;Integrated Security=True;Encrypt=True;TrustServerCertificate=False;"
  },
  "DeveloperMode": {
    "EnableFormLogin": false,
    "EnableUserHeaderAuth": false,
    "EnableDemoSeedData": false
  },
  "AllowedHosts": "*"
}
\end{lstlisting}

Press \textbf{Ctrl + S} to save the file. That is the entire file.

\bigskip
\textbf{What each part means:}

\begin{itemize}
  \item \texttt{Server=clipvwbpod02} ---
        The hostname of the production SQL Server.
        Do not change this unless a DBA tells you to.
  \item \texttt{Database=ClientsAppAttestation} ---
        The database name.
        Do not change this.
  \item \texttt{Integrated Security=True} ---
        Connect to SQL Server using the Windows identity of the IIS application pool.
        No username or password is stored in this file.
  \item \texttt{Encrypt=True} ---
        All traffic between the app and the database is encrypted.
        Required in production.
  \item \texttt{TrustServerCertificate=False} ---
        The SQL Server's certificate must be valid.
        Do not change this to \texttt{True} in production.
  \item \texttt{EnableFormLogin: false} ---
        Disables the developer-only form login endpoint.
        Must be \texttt{false}.
  \item \texttt{EnableUserHeaderAuth: false} ---
        Disables the \texttt{X-User-Id} header shortcut that bypasses Windows auth.
        Must be \texttt{false}.
  \item \texttt{EnableDemoSeedData: false} ---
        Prevents fictional test users from being inserted into the database.
        Must be \texttt{false}.
\end{itemize}

\begin{warnbox}
\textbf{Never change any DeveloperMode value to \texttt{true} in production.}
\texttt{EnableFormLogin} and \texttt{EnableUserHeaderAuth} bypass Windows Authentication
and allow any caller to impersonate any user without a password.
\texttt{EnableDemoSeedData} will insert Bruce Banner, Diana Prince, Selina Kyle, and
$\approx$625 fictional test records into the live database.
\end{warnbox}

% ================================================================
\section{If I Want to Set the Connection String in IIS Instead}
\label{sec:iisonly}
% ================================================================

If your deployment pipeline overwrites \texttt{appsettings.Production.json}
on every deploy, set the connection string as an environment variable in IIS.
Environment variables always override the settings file.

\subsection*{Method A --- IIS Manager (no file editing needed)}

\begin{enumerate}
  \item Press \textbf{Win + R}, type \texttt{inetmgr}, press \textbf{Enter}.
        IIS Manager opens.
  \item In the left panel, expand the server name, expand \textbf{Sites},
        and click on \textbf{DashyDashboard}.
  \item In the centre panel, double-click \textbf{Configuration Editor}.
  \item At the top of the Configuration Editor, find the \textbf{Section} dropdown.
        Click it and navigate to \texttt{system.webServer} $\rightarrow$ \texttt{aspNetCore}.
  \item In the table that appears in the centre, find the row labelled
        \textbf{environmentVariables}.
        Click the grey \textbf{[\ldots]} button at the right end of that row.
        A Collection Editor window opens.
  \item In the top-right of the Collection Editor, click \textbf{Add Item}.
        A new blank row appears in the list.
  \item Click the \textbf{name} cell of that new row and type:
\begin{lstlisting}
ConnectionStrings__Default
\end{lstlisting}
        (Two underscores between \texttt{ConnectionStrings} and \texttt{Default}.)
  \item Click the \textbf{value} cell and type:
\begin{lstlisting}
Server=clipvwbpod02;Database=ClientsAppAttestation;Integrated Security=True;Encrypt=True;TrustServerCertificate=False;
\end{lstlisting}
  \item Click \textbf{Apply} in the top-right Actions panel. Close the Collection Editor.
  \item Click \textbf{Apply} again in the main Configuration Editor panel.
  \item In the left panel, click \textbf{Application Pools}.
        Right-click \textbf{DashyDashboard}. Click \textbf{Recycle}.
\end{enumerate}

\subsection*{Method B --- Edit web.config directly}

Open \texttt{C:\textbackslash inetpub\textbackslash DashyDashboard\textbackslash web.config}
in Notepad.
Find the \texttt{<aspNetCore>} element.
Add both variables inside an \texttt{<environmentVariables>} block:

\begin{lstlisting}[language={}]
<aspNetCore processPath=".\DashyDashboard.Api.exe"
            stdoutLogEnabled="false"
            stdoutLogFile=".\logs\stdout"
            hostingModel="inprocess">
  <environmentVariables>
    <environmentVariable
        name="ASPNETCORE_ENVIRONMENT"
        value="Production" />
    <environmentVariable
        name="ConnectionStrings__Default"
        value="Server=clipvwbpod02;Database=ClientsAppAttestation;Integrated Security=True;Encrypt=True;TrustServerCertificate=False;" />
  </environmentVariables>
</aspNetCore>
\end{lstlisting}

Press \textbf{Ctrl + S}. No IIS restart needed.
IIS applies \texttt{web.config} changes immediately.

\begin{infobox}
\textbf{Why two underscores?}
ASP.NET Core converts environment variable names to configuration keys by replacing
\texttt{\_\_} (double underscore) with \texttt{:} (colon).
So the variable \texttt{ConnectionStrings\_\_Default} maps to the key
\texttt{ConnectionStrings:Default}, which is exactly the key the app reads.
\end{infobox}

% ================================================================
\section{Windows Authentication Setup}
\label{sec:winauth}
% ================================================================

The app uses Windows/Negotiate authentication.
IIS identifies the user from their Windows login and passes that identity to the app.
You must enable Windows Authentication and disable Anonymous Authentication.

\subsection*{Step 1 --- Enable the Windows Authentication feature in Windows Server}

Do this once. If Windows Authentication is already listed in IIS, skip to Step~2.

\begin{enumerate}
  \item Click \textbf{Start}. Type \texttt{Server Manager}. Press \textbf{Enter}.
  \item Click \textbf{Manage} in the top-right corner.
        Click \textbf{Add Roles and Features}.
  \item Click \textbf{Next} through the first three pages
        (Before You Begin, Installation Type, Server Selection).
  \item On \textbf{Server Roles}, find and expand:
        \textbf{Web Server (IIS)} $\rightarrow$ \textbf{Web Server} $\rightarrow$ \textbf{Security}.
  \item Tick the checkbox next to \textbf{Windows Authentication}.
        If a dialog says ``Add features that are required?'', click \textbf{Add Features}.
  \item Click \textbf{Next} twice more, then click \textbf{Install}.
  \item Wait for the installation bar to finish. Click \textbf{Close}.
\end{enumerate}

\subsection*{Step 2 --- Create the application pool}

\begin{enumerate}
  \item Press \textbf{Win + R}. Type \texttt{inetmgr}. Press \textbf{Enter}.
  \item In the left panel, click \textbf{Application Pools}.
  \item In the right-side Actions panel, click \textbf{Add Application Pool\ldots}
  \item In the dialog that appears, fill in:
    \begin{itemize}
      \item \textbf{Name:} \quad \texttt{DashyDashboard}
      \item \textbf{.NET CLR version:} \quad \texttt{No Managed Code}
            (this is correct --- the app is self-contained)
      \item \textbf{Managed pipeline mode:} \quad \texttt{Integrated}
    \end{itemize}
  \item Click \textbf{OK}.
\end{enumerate}

\subsection*{Step 3 --- Create the website}

\begin{enumerate}
  \item In the left panel, right-click \textbf{Sites} and click \textbf{Add Website\ldots}
  \item Fill in every field:
    \begin{itemize}
      \item \textbf{Site name:} \quad \texttt{DashyDashboard}
      \item \textbf{Application pool:} \quad click \textbf{Select\ldots} and pick
            \texttt{DashyDashboard}
      \item \textbf{Physical path:} \quad the folder on this server where you copied
            the published output, for example \texttt{C:\textbackslash inetpub\textbackslash DashyDashboard}
      \item \textbf{Binding type:} \quad \texttt{https}
      \item \textbf{IP address:} \quad \texttt{All Unassigned}
      \item \textbf{Port:} \quad \texttt{443}
      \item \textbf{Host name:} \quad the DNS name users will type in their browser,
            for example \texttt{dashy.broadridge.internal}
      \item \textbf{SSL certificate:} \quad pick the certificate that matches
            your host name from the dropdown
    \end{itemize}
  \item Click \textbf{OK}.
\end{enumerate}

\begin{infobox}
The application server hostname (what users type in the browser) is not stored anywhere
in the repository. It is whatever DNS name you enter in the Host name field above.
Tell users to browse to that name.
\end{infobox}

\subsection*{Step 4 --- Configure authentication}

\begin{enumerate}
  \item In the left panel, click on your \textbf{DashyDashboard} site.
  \item In the centre panel, double-click \textbf{Authentication}.
  \item Click \textbf{Anonymous Authentication} once (just to select it).
        In the right Actions panel, click \textbf{Disable}.
  \item Click \textbf{Windows Authentication} once to select it.
        In the right Actions panel, click \textbf{Enable}.
  \item With \textbf{Windows Authentication} still selected, click
        \textbf{Providers\ldots} in the right Actions panel.
        A small dialog opens showing a list.
  \item The list must show \textbf{Negotiate} above \textbf{NTLM}.
        If it does not, click \textbf{Add}, select \texttt{Negotiate}, click \textbf{OK},
        then use the \textbf{Move Up} button to put Negotiate above NTLM.
  \item Click \textbf{OK} to close the Providers dialog.
\end{enumerate}

The authentication panel should now show:

\bigskip
\begin{tabular}{@{}ll@{}}
\toprule
Authentication type & Status \\
\midrule
Anonymous Authentication & Disabled \\
Windows Authentication   & Enabled \\
\bottomrule
\end{tabular}

\subsection*{Step 5 --- Grant the app pool identity access to SQL Server}

The app connects to SQL Server using the Windows account of the IIS application pool.
Ask the DBA to run this SQL on \texttt{clipvwbpod02}:

\begin{lstlisting}[language=SQL]
USE ClientsAppAttestation;

CREATE LOGIN [IIS APPPOOL\DashyDashboard]
    FROM WINDOWS
    WITH DEFAULT_DATABASE = ClientsAppAttestation;

CREATE USER [IIS APPPOOL\DashyDashboard]
    FOR LOGIN [IIS APPPOOL\DashyDashboard];

ALTER ROLE db_datareader
    ADD MEMBER [IIS APPPOOL\DashyDashboard];

ALTER ROLE db_datawriter
    ADD MEMBER [IIS APPPOOL\DashyDashboard];
\end{lstlisting}

If a domain service account is used instead of the default pool identity,
replace every occurrence of
\texttt{IIS APPPOOL\textbackslash DashyDashboard}
with that account name.

% ================================================================
\section{Build and Publish}
\label{sec:publish}
% ================================================================

All commands in this section are run on the \textbf{build machine}, not the IIS server.

\subsection*{Step 1 --- Open a Command Prompt at the repo root}

\begin{enumerate}
  \item Press \textbf{Win + R}. Type \texttt{cmd}. Press \textbf{Enter}.
  \item Type \texttt{cd} followed by a space, then the full path to the repository root.
        Press \textbf{Enter}. For example:
\begin{lstlisting}
cd C:\projects\dashythedashdashboard
\end{lstlisting}
        Replace \texttt{C:\textbackslash projects\textbackslash dashythedashdashboard}
        with wherever the repository actually lives on your machine.
\end{enumerate}

\subsection*{Step 2 --- Confirm prerequisites}

\begin{lstlisting}
dotnet --version
\end{lstlisting}
Output must begin with \texttt{7.} or higher (e.g.\ \texttt{7.0.407} or \texttt{8.0.421}).

\begin{lstlisting}
node --version
\end{lstlisting}
Output must begin with \texttt{v18.} or higher.

If either is wrong or missing, go back to Section~\ref{sec:prereqs} and install.

\subsection*{Step 3 --- Edit the connection string (if you have not already)}

Follow Sections~\ref{sec:whichfile} and~\ref{sec:whatpaste}.
Save the file before continuing.

\subsection*{Step 4 --- Run the publish command}

Type this command exactly and press \textbf{Enter}:

\begin{lstlisting}
dotnet publish DashyDashboard/DashyDashboard.Api/DashyDashboard.Api.csproj -c Release -o publish/DashyDashboard.Api
\end{lstlisting}

You do not need to run any npm or frontend commands manually.
The publish step does all of the following automatically:

\begin{enumerate}
  \item Runs \texttt{npm ci} inside \texttt{DashyDashboard/DashyDashboard.Frontend}
        to install all Node packages.
  \item Runs \texttt{npm run build} to compile the React app into static files.
  \item Copies those static files into the API's \texttt{wwwroot} folder.
  \item Compiles the .NET API in Release mode.
  \item Copies every file needed to run the application into
        \texttt{publish\textbackslash DashyDashboard.Api}.
\end{enumerate}

When it finishes successfully, the last line of output will be:
\begin{lstlisting}
Build succeeded.
\end{lstlisting}

If you see an error about \texttt{npm} not being found, install Node.js 18
(Section~\ref{sec:prereqs}) and try again.

\subsection*{Step 5 --- Copy the output to the IIS server}

The published output is now in the folder
\texttt{publish\textbackslash DashyDashboard.Api} on the build machine.
Copy everything inside that folder to the site physical path on the IIS server.

Run this command from the repository root (replace \texttt{your-iis-server} with
the actual server name or IP address):

\begin{lstlisting}
robocopy publish\DashyDashboard.Api \\your-iis-server\c$\inetpub\DashyDashboard /MIR /Z
\end{lstlisting}

Alternatively, copy the folder manually using Windows File Explorer over a network share.

\subsection*{Step 6 --- Run database migrations}

\begin{warnbox}
\textbf{Do not skip this.}
In production, the app does not apply migrations automatically at startup.
If you skip this step, the tables will not exist and every API call will fail.
\end{warnbox}

Run this from a machine that:
(a) has the .NET 7 SDK (or later) installed, and
(b) can reach \texttt{clipvwbpod02} on the network
(this can be the build machine if it has network access to the database server).

In the same Command Prompt, navigate to the API project:

\begin{lstlisting}
cd DashyDashboard/DashyDashboard.Api
\end{lstlisting}

Check that the EF Core CLI is installed:

\begin{lstlisting}
dotnet ef --version
\end{lstlisting}

If that command prints ``command not found'' or fails, install it:

\begin{lstlisting}
dotnet tool install --global dotnet-ef
\end{lstlisting}

Then run the migrations:

\begin{lstlisting}
dotnet ef database update --connection "Server=clipvwbpod02;Database=ClientsAppAttestation;Integrated Security=True;Encrypt=True;TrustServerCertificate=False;"
\end{lstlisting}

You will see each migration name printed one by one.
The last line will be:
\begin{lstlisting}
Done.
\end{lstlisting}

If you get an SSL certificate error, ask the DBA to confirm that the SQL Server
certificate on \texttt{clipvwbpod02} is valid and trusted by the machine you are
running from. If the DBA confirms the cert is internal-only,
change \texttt{TrustServerCertificate=False} to \texttt{True}
\textbf{only for this migration command}, not in the deployed settings file.

% ================================================================
\section{How to Keep Demo Data Out of Production}
\label{sec:demodata}
% ================================================================

The application includes seed data with fictional users:
\textbf{Bruce Banner, Diana Prince, Selina Kyle}, and approximately 625 additional
test records.

The seed code runs \textbf{only when both} of these conditions are true simultaneously:

\begin{itemize}
  \item \texttt{ASPNETCORE\_ENVIRONMENT} equals \texttt{Development}, \textbf{and}
  \item \texttt{DeveloperMode:EnableDemoSeedData} equals \texttt{true}
\end{itemize}

In production, neither condition is true if you follow this guide.
The \texttt{appsettings.Production.json} from Section~\ref{sec:whatpaste} already
sets \texttt{EnableDemoSeedData} to \texttt{false},
and \texttt{ASPNETCORE\_ENVIRONMENT} is \texttt{Production}.

\medskip
\textbf{Verification checklist --- check each of these before go-live:}

\begin{enumerate}
  \item Open \texttt{C:\textbackslash inetpub\textbackslash DashyDashboard\textbackslash appsettings.Production.json}
        in Notepad.
        Confirm the \texttt{DeveloperMode} block reads exactly:
\begin{lstlisting}
"DeveloperMode": {
  "EnableFormLogin": false,
  "EnableUserHeaderAuth": false,
  "EnableDemoSeedData": false
}
\end{lstlisting}
        All three values must be \texttt{false} (no quotes around it,
        not \texttt{"false"}, not missing).
  \item Open IIS Manager. Click on the site.
        Double-click \textbf{Configuration Editor}.
        Navigate to \texttt{system.webServer/aspNetCore}.
        Click \textbf{[\ldots]} next to \textbf{environmentVariables}.
        Confirm there is no variable named
        \texttt{DeveloperMode\_\_EnableDemoSeedData} with value \texttt{true}.
        Environment variables override the file.
  \item Confirm the variable \texttt{ASPNETCORE\_ENVIRONMENT} is set to
        \texttt{Production} (not \texttt{Development}) in either \texttt{web.config}
        or the IIS Configuration Editor.
\end{enumerate}

% ================================================================
\section{What Remains Before First Production Run}
\label{sec:blockers}
% ================================================================

\begin{infobox}
\textbf{Schema alignment is complete.}
All blockers from the original pre-development audit have been resolved in the
current codebase:
\texttt{AssociateID} is \texttt{varchar(50)} throughout;
\texttt{ToolID} is \texttt{int} throughout;
all ten spec tables exist (\texttt{SuperUsers}, \texttt{IFHGFHMapping},
\texttt{AttestationLogs}, \texttt{LoginLogs}, and the rest);
login attempts are written to \texttt{LoginLogs} on every request.
The build is clean (0 errors).
\end{infobox}

What remains is operational, not code.

\subsection*{Mandatory --- Run EF Core migrations}

Covered in Section~\ref{sec:publish}, Step~6.
The migrations create all ten tables.
The database starts empty after migration --- no application data is inserted.

\subsection*{Mandatory --- Import production data from Excel}

The application tables are populated from the authoritative Excel files maintained
by the operations team.
The source files referenced in the schema documentation are:

\begin{itemize}
  \item \texttt{C:\textbackslash Prakash Files\textbackslash Dash\textbackslash Users.xlsx}
        $\rightarrow$ \texttt{dbo.Users}
  \item \texttt{C:\textbackslash Prakash Files\textbackslash Dash\textbackslash Clients.xlsx}
        $\rightarrow$ \texttt{dbo.Clients}
  \item \texttt{C:\textbackslash Prakash Files\textbackslash Dash\textbackslash ClientTools.xlsx}
        $\rightarrow$ \texttt{dbo.ClientTools}
  \item \texttt{C:\textbackslash Prakash Files\textbackslash Dash\textbackslash UsersToolAccess.xlsx}
        $\rightarrow$ \texttt{dbo.UsersToolAccess}
  \item \texttt{C:\textbackslash Prakash Files\textbackslash Dash\textbackslash Department IFH GFH Mapping.xlsx}
        $\rightarrow$ \texttt{dbo.IFHGFHMapping}
\end{itemize}

Use the \texttt{OPENROWSET} / ACE OLEDB import pattern from Appendix~A of the
schema PDF (\texttt{ClientsAppAttestation\_Table\_Structure\_Documentation.pdf}).
Enable Ad Hoc Distributed Queries on the SQL Server first:

\begin{lstlisting}[language=SQL]
EXEC sp_configure 'show advanced options', 1; RECONFIGURE;
EXEC sp_configure 'Ad Hoc Distributed Queries', 1; RECONFIGURE;
\end{lstlisting}

The SQL Server service account on \texttt{clipvwbpod02} must be able to read the
file path above (either a local path on the SQL Server machine, or a UNC path
accessible from it).

\subsection*{Mandatory --- Populate SuperUsers for the go-live team}

Once Users are imported, insert a row in \texttt{dbo.SuperUsers}
for each person who needs Admin, VP, GFH, or IFH access.
Ask the DBA to run:

\begin{lstlisting}[language=SQL]
USE ClientsAppAttestation;

INSERT INTO dbo.SuperUsers
    (AssociateID, RoleName, Department, AccessLevel, IsActive, CreatedBy)
VALUES
    ('173926', 'Admin', 'DTC Settlements', 'Full', 'TRUE', 'System'),
    ('179335', 'VP',    'International',   'Approver', 'TRUE', 'System');
-- Add one row per admin/VP/GFH/IFH user.
\end{lstlisting}

Use the real \texttt{AssociateID} values (VARCHAR, e.g.\ \texttt{'173926'}) and
the real \texttt{RoleName} values: \texttt{Admin}, \texttt{VP}, \texttt{GFH},
or \texttt{IFH}.

\subsection*{Critical --- Set the UserName column for Windows Authentication}

This step is the most common reason users reach the app but see nothing.

When a user opens the app in their browser, IIS identifies them via Windows
Authentication (Kerberos/NTLM) and passes their identity to the application as
\texttt{HttpContext.User.Identity.Name}.
The app then queries:

\begin{lstlisting}[language=SQL]
SELECT * FROM dbo.Users WHERE UserName = '<value IIS sent>';
\end{lstlisting}

If no row matches, the user is unresolved and the app behaves as if nobody is
logged in.

\textbf{What value does IIS send?}
On most corporate Windows setups, \texttt{User.Identity.Name} is
\texttt{DOMAIN\textbackslash loginname}, for example \texttt{BROADRIDGE\textbackslash KalapurayilP}.
The \texttt{UserName} column in \texttt{dbo.Users} must contain exactly
that string for each user.

Verify what IIS sends on your network by enabling stdout logging temporarily
(see the Troubleshooting Guide, Section~2) and looking for the identity in the
application logs, or ask the infrastructure team what domain prefix IIS uses.

Once confirmed, update the column in bulk:

\begin{lstlisting}[language=SQL]
-- Example: if IIS sends BROADRIDGE\<loginname>
UPDATE dbo.Users
SET UserName = 'BROADRIDGE\' + UserName
WHERE UserName NOT LIKE 'BROADRIDGE\%';
\end{lstlisting}

If IIS sends the bare login name without a domain prefix, the UserName values
already imported from Excel (e.g.\ \texttt{KalapurayilP}) will match directly
and no update is needed.

\subsection*{Remaining non-blocking limitations}

\begin{itemize}
  \item \textbf{No mobile layout.}
        The frontend targets desktop viewports (1280\,px and wider).
  \item \textbf{Dev form login accepts any password.}
        The \texttt{POST /api/auth/login} endpoint is disabled in Production
        (\texttt{EnableFormLogin: false}) but in Development it accepts any
        non-empty password without verification.
        This is intentional --- Production never uses form login.
  \item \textbf{ClientTools and UsersToolAccess use the import-table structure.}
        \texttt{dbo.ClientTools} uses \texttt{UserID int identity} as its PK
        and \texttt{[Application Name] varchar(255)} for the tool name, matching
        the flat Excel import design.
        The normalised design (composite key on ClientID+ToolName) is enforced
        as a unique constraint but not the primary key.
        This is intentional and matches the production database.
\end{itemize}

\end{document}
